\section*{Contoh}
\addcontentsline{toc}{section}{Contoh}

\begin{contoh}[Menghitung FIRST]
Untuk grammar
\[
E \to T E', \quad E' \to + T E' \mid \varepsilon, \quad
T \to F T', \quad T' \to * F T' \mid \varepsilon, \quad
F \to (E) \mid \code{id},
\]
FIRST(\(F\)) = \(\{\code{id}, (\}\). Karena \(T \to F T'\) dan
\(F\) tidak nullable, FIRST(\(T\)) = FIRST(\(F\)). Demikian pula
FIRST(\(E\)) = FIRST(\(T\)). Untuk \(E'\), FIRST(\(E'\)) =
\(\{+, \varepsilon\}\).
\end{contoh}

\begin{contoh}[Menghitung FOLLOW]
FOLLOW(\(E\)) = \(\{\$, )\}\) karena \(E\) adalah simbol awal dan
muncul di dalam kurung. FOLLOW(\(E'\)) sama dengan FOLLOW(\(E\)).
FOLLOW(\(T\)) = \(\{\$, ), +\}\) karena \(T\) diikuti \(E'\), dan
FIRST(\(E'\)) memuat \(+\). FOLLOW(\(F\)) = \(\{\$, ), *, +\}\)
karena \(F\) diikuti \(T'\), dan FIRST(\(T'\)) memuat \(*\).
\end{contoh}

\begin{contoh}[Mengisi tabel LL(1)]
Aturan \(E \to T E'\) diisikan pada kolom FIRST(\(T\)) =
\(\{\code{id}, (\}\). Aturan \(E' \to \varepsilon\) diisikan pada
kolom FOLLOW(\(E'\)) = \(\{\$, )\}\). Aturan
\(T' \to * F T'\) diisikan pada kolom \(*\). Sel yang tidak terisi
menandakan kesalahan sintaksis.
\end{contoh}

\begin{contoh}[Penelusuran stack]
Untuk masukan \(\code{id} + \code{id}\ \$\), stack dimulai dengan
\(\$E\). Aturan \(E \to TE'\) menghasilkan stack \(\$E'T\).
Aturan \(T \to FT'\) menghasilkan \(\$E'T'F\). Aturan
\(F \to \code{id}\) menghasilkan \(\$E'T'\code{id}\), lalu
\code{id} dicocokkan dan di-pop. Setelah \(T' \to \varepsilon\),
\(E' \to +TE'\), dan seterusnya, parser menerima masukan ketika
stack hanya berisi \(\$\).
\end{contoh}

\begin{contoh}[Pemulihan titik koma]
Untuk masukan
\begin{lstlisting}[language=C]
int a = 1
return 0;
\end{lstlisting}
parser melaporkan \code{baris 1, kolom 10: diharapkan ';' setelah
deklarasi, ditemukan 'return'}. Setelah itu, parser menyinkronkan
ke titik koma berikutnya dan melanjutkan. Program tetap dapat
dianalisis lebih lanjut meskipun mengandung kesalahan.
\end{contoh}
